\documentclass[10.5pt, a4paper]{article}

\usepackage[T1]{fontenc}
\usepackage[utf8]{inputenc}
\usepackage{tgpagella}
\usepackage{tgcursor}
\usepackage{microtype}
\usepackage[spanish, es-noshorthands]{babel}
\usepackage[table, xcdraw]{xcolor}
\usepackage{booktabs}
\usepackage{tabularx}
\usepackage{array}
\usepackage{amsmath, amssymb}
\usepackage[most]{tcolorbox}
\usepackage{listings}
\usepackage[margin=1.9cm]{geometry}
\usepackage{fancyhdr}

\lstset{
    basicstyle=\ttfamily\footnotesize,
    backgroundcolor=\color{gray!5},
    frame=single,
    rulecolor=\color{gray!40},
    showstringspaces=false,
    breaklines=true,
    breakatwhitespace=false,
    columns=flexible,
    literate=
      {á}{{\'a}}1 {é}{{\'e}}1 {í}{{\'i}}1 {ó}{{\'o}}1 {ú}{{\'u}}1
      {Á}{{\'A}}1 {É}{{\'E}}1 {Í}{{\'I}}1 {Ó}{{\'O}}1 {Ú}{{\'U}}1
      {ñ}{{\~n}}1 {Ñ}{{\~N}}1
}

\pagestyle{fancy}
\fancyhf{}
\setlength{\headheight}{16pt}
\lhead{\textbf{IMT-342} $\mid$ Robótica}
\chead{Esquema de Resultados Reproducibles — Hito 3}
\rhead{Contratos de Datos y Código Offline}
\lfoot{Evidencia Trazable y Reproducible}
\rfoot{Página \thepage}
\renewcommand{\headrulewidth}{0.4pt}
\renewcommand{\footrulewidth}{0.4pt}

\definecolor{ucbblue}{RGB}{0, 51, 102}

\begin{document}

\begin{center}
    \Large\textbf{\textcolor{ucbblue}{Esquema de Resultados Reproducibles, Contratos CSV y Referencia Offline (Hito 3)}}
\end{center}

\section{Estructura Propuesta del Repositorio de Evidencias}
Todo equipo debe organizar sus datos brutos, resultados, scripts y aprobaciones reales bajo una jerarquía reproducible. Si una magnitud no existe o no fue ejecutada, registrar \texttt{NA} junto con su razón en la bitácora; \textbf{nunca sustituir datos ausentes por cero}.

\begin{lstlisting}
hito3_evidence/
  README.md                   # Pasos exactos de ejecucion y versiones de entorno
  manifest.yaml               # Metadatos obligatorios del ensayo (ver Seccion 3)
  tests/
    A01_inventory.csv         # Inventario de controlador, RobotWare, ROS, URDF
    B01_joint_states.csv      # Lecturas pasivas de 6 juntas (3 estados estacionarios)
    C01_fk.csv                # Comparaciones FK (ep, eR y casos sinteticos de diagnostico)
    E01_ik.csv                # Ramas IK candidatas, limites y residuos FK
    F01_jacobian.csv          # Errores eJv, eJw, convergencia FD y SVD escalado
    G01_lspb.csv              # Perfil LSPB offline, fronteras y extremos analiticos
    D00_command_review.csv    # Revision de contrato de interfaz y acuse/rechazo sin mov.
    D01_tracking.csv          # SOLAMENTE si hubo ejecucion fisica autorizada (GO)
  scripts/                    # Codigos reproducibles versionados (offline)
  logs/                       # Archivos brutos de consola, red y eventos de error
  figures/                    # Graficas generadas desde los CSV anteriores
  approvals/                  # Exclusivamente autorizaciones GO emitidas realmente
\end{lstlisting}

\section{Esquemas CSV Expandidos sin Ambigüedad}
Los vectores no deben empaquetarse en una sola celda con corchetes; en CSV deben expandirse explícitamente a columnas escalares numeradas. Las marcas temporales deben escribirse en formato \texttt{ISO8601} con desplazamiento horario o como tiempo monotónico declarando dominio y origen. Las orientaciones deben declararse en matriz de rotación o cuaternión especificando el orden (\texttt{xyzw} vs. \texttt{wxyz}). No inventar \texttt{source\_feedback} si se desconoce.

\subsection{Esquema \texttt{B01\_joint\_states.csv} (Estado Articular — Etapa B)}
\begin{lstlisting}
test_id,source,sample_index,timestamp_source,clock_domain,timestamp_receive,joint_name_1,joint_name_2,joint_name_3,joint_name_4,joint_name_5,joint_name_6,q_actual_1,q_actual_2,q_actual_3,q_actual_4,q_actual_5,q_actual_6,units,controller_state,flags,provenance_file
\end{lstlisting}

\subsection{Esquema \texttt{D01\_tracking.csv} (Consigna y Seguimiento — Solo si Autorizado)}
\begin{lstlisting}
test_id,timestamp,clock_domain,q_ref_1,q_ref_2,q_ref_3,q_ref_4,q_ref_5,q_ref_6,q_actual_1,q_actual_2,q_actual_3,q_actual_4,q_actual_5,q_actual_6,pose_ref_px,pose_ref_py,pose_ref_pz,pose_ref_rot_repr,pose_actual_px,pose_actual_py,pose_actual_pz,pose_actual_rot_repr,base_frame,tcp_id,source_command,source_feedback,ack_status,execution_status,safety_go_id,flags,pos_units,ang_units
\end{lstlisting}

\section{Plantilla Mínima de \texttt{manifest.yaml}}
\begin{lstlisting}
project: IMT-342_Hito3
team: null
test_id: null
date: null                     # ISO8601
robot_controller_robotware:
  robot_model: ABB_IRB120
  controller_model: null       # ej. IRC5 (verificar localmente)
  robotware_version: null
  options_verified: []
ros_driver_commit:
  ros_distro: null
  driver_repo_commit: null
urdf_hash: null
dh_parameters_version: null
frames:
  base_frame: null
  flange_frame: null
  tcp_frame: null
joints:
  joint_order: []
  signs: []
  offsets: []
units:
  position: m
  angle: rad
  angular_velocity: rad/s
timebase:
  clock_domain: null
  synchronization_method: null
reference_source_independence:
  reference_name: null
  independent_of_own_urdf: null
thresholds_version_approval:
  criteria_version: null
  approval_status: pending
safety_status: NO-GO
data_file_hashes: {}
operator_reviewer:
  operator: null
  reviewer: null
known_limitations: []
\end{lstlisting}

\section{Código Didáctico de Referencia Offline (Modelo Sintético 2R — No ABB)}
\begin{tcolorbox}[colback=yellow!10, colframe=orange!85!black, left=2mm, right=2mm, top=1mm, bottom=1mm]
    \footnotesize\raggedright
    \textbf{Alcance estricto:} Este script utiliza Python 3 y \texttt{numpy} exclusivamente para referencia didáctica offline sobre un mecanismo planar 2R sintético. No implementa APIs de ROS ni control físico, ni constituye un validador final del ABB IRB120.
\end{tcolorbox}

\begin{lstlisting}[language=Python]
import numpy as np

a1, a2 = 1.0, 0.5  # modelo sintetico, NO ABB
q = np.array([np.pi / 6, -np.pi / 4])
qd = np.array([0.1, -0.2])

def fk_2r(x):
    u, v = x
    return np.array([a1*np.cos(u) + a2*np.cos(u + v),
                     a1*np.sin(u) + a2*np.sin(u + v)])

def jac_2r(x):
    u, v = x
    return np.array([
        [-a1*np.sin(u) - a2*np.sin(u + v), -a2*np.sin(u + v)],
        [ a1*np.cos(u) + a2*np.cos(u + v),  a2*np.cos(u + v)],
    ])

J = jac_2r(q)
print("p:", fk_2r(q))
print("J:", J)
print("v:", J @ qd)

for h in (1e-2, 1e-3, 1e-4, 1e-5, 1e-6):
    eye = np.eye(2)
    J_fd = np.column_stack([
        (fk_2r(q + h*eye[j]) - fk_2r(q - h*eye[j])) / (2*h)
        for j in range(2)
    ])
    print("h:", h, "eJv_F:", np.linalg.norm(J - J_fd, ord="fro"))
\end{lstlisting}

\noindent\textbf{Validación analítica esperada:} $p \approx [1.34898832,\; 0.37059048]^T\;\mathrm{m}$, $J \approx \begin{bmatrix} -0.37059048 & 0.12940952 \\ 1.34898832 & 0.48296291 \end{bmatrix}\;\mathrm{m/rad}$, $v \approx [-0.06294095,\; 0.03830625]^T\;\mathrm{m/s}$. Los errores numéricos \texttt{eJv\_F} deben ejecutarse y documentarse en el entorno real del estudiante con versión de Python/\texttt{numpy} registrada en bitácora.

\end{document}
